
\renewcommand{\footerfilename}{identify.tex}

\section{IRAF identify task}

Identify, well described by its help file, has few tricks of use for spectroscopy.

It can create a WCS using either emissin lines (in the case of a calibration lamp)
or it may be drawn from the image's absorption features directly.

The key is to fine tune the parameters (epar/lpar etc) for the instrument. A bit of
work once, perhaps a few fine tweaks for a bad night.

Create a line list that matches your instrument. This can be done with a long list,
an initial (tedious) extraction. Then the lines that were actually used are listed
in the database entry for that file. Take that list and make a new list to use going
forward. Note: The list changes when the bulb changes -- usually with night-to-night
power supply fluxuations or when it is physically swapped out.

The short list makes auto line matching much faster.


Initial Conditions:

\vspace{-.15cm}
\begin{enumerate}\addtolength{\itemsep}{-0.5\baselineskip}
%\setcounter{enumi}{N}
   \item   All files related to a calibration lamp image should start with the same starting
   $(X<Y)$ coordinates.
   \item   extract a trace from each file of the science sequence. This produces a 1D-fits image.
   \item   apply identify to the middle file of that sequence.
   \item   This produces a mapping of $(X) \rightarrow (lambda)$. 
   \item   Use this identify file, with re-identify for the other spectra of the sequence. The
     assuption is that the shifts are small -- they can be constrained with the lparams.
\end{enumerate}



Start with the main trick of trimming all the spectra for a camera to tight blue
and red extents.  This helps to avoid Runge phenomenon with polynomial
fitting -- trim to exclude $X$ values that carry no information.

A critical thing to remember is that a astronomical spectrum represents flux generated
by discrete number of elements (and a few other aspects, like temperature, the local
gravity field, the B field, etc) that are combined together. One region of the
spectrum is not necessarily informed by any adjacent (different element) information.
So, using mathematics that starts with phrases like ``given a continuous function, ...''
we run into the problem that each element spans a finite distance in the independent
axis. 

All this translates into adding noise to the very quantities we want to describe.

\subsection{Get the linelist right} \ref{identify_linelist}

The linelist \index{identify!linelist} contains a list of the line position
in angstroms together with the name of the line. Do a hand-fit of the lamp
the old fashioned way, guess/fit etc. Using the \llbox{D} key to delete
poorly matched lines, save and exit the function. Under the \dhl{database}
location, find the file you just created. Within that file is the list of
lines, their angstrom values, and a confidence of the fit. Using well fit
lines, develop the \dhl{practical line list}.

A decent, concise but accurate, list of line that can be practically fit,
the \dhl{reidentify} step will be successful. 

Prepare a hydrogen list to be used with absorption features.

\subsection{Trick - Calibrate a calibration image}

In order to extract a trace; the trace needs to be brighter than the sky and
have a gaussian cross-section. A python program \index{python!fakedispersion}
will draw a gaussian across the calibration image. You can then use apall to
extract that trace. At this point you need to use the identify task to wavelength
calibrate the file. The goal is to produce a handy plot for later reference.
One easier way to do this is to use an A0V taken within the same context of
the cal lamp image. Then use the absorption line method and the hydrogen.dat
file to get a good guess of the WCS. Use nhedit to add that to the faketrace,
then use the reidentify task together with the emission.dat file to get
a closer approximation. Add the labels to the file, save and print that
file. This serves as a way to disamabiguate the small features in the blue
from Argon lines -- applied to the intensity profile from the specific cal lamp.


\subsection{Trick - print the plot of well fitted lines}

During the identify step, the \dhl{user coordinates} command will overlay
prose onto the graph. This makes identification by hand easier later. 




\ltodo{identify}{Fill in more tricks.}

\begingroup \fontsize{10pt}{10pt}
\selectfont
%%\begin{Verbatim} [commandchars=\\\{\}]
\begin{verbatim} 
With Identify - use lables to denote the lines.
do a snap
screencapture the image
gimp make the bg transparent, and the colors agree with seeing the lines.
Then add this as launcher:


/usr/bin/pqiv -c  /home/wayne/Pictures/TonyCalsAlphaWhite.png
\end{verbatim}
\endgroup
%% \end{Verbatim}

\subsection{Fine Details}

If you trim your images to the same starting pixel/wavelength you can specify that starting
position with the crval and cdelt values. These can be observed in the header of a final
output file and fed-back into this mix.

Traces are usually best in the middle of the spectrum -- where the brightest features are
located. 

database is where the answers are stored. There naming convention is pretty consistent.

